\chapter{Eksperymenty i analiza wyników}
\label{ch:eksperymenty}

\section{Opis zbioru danych wykorzystanego w badaniach}
\label{sec:zbior-danych}

% TODO: Opisać użyty korpus mowy, np. LibriSpeech train-clean-100 lub train-clean-360, sposób podziału danych,
% przygotowanie próbek, długość segmentów i ograniczenia wynikające z dostępności danych.

\section{Konfiguracja eksperymentów}
\label{sec:konfiguracja-eksperymentow}

% TODO: Udokumentować konfiguracje treningu: liczbę kroków, batch, akumulację gradientu, architekturę,
% parametry RVQ, funkcje kosztu, checkpointy i sprzęt. Warto odwołać się do plików w katalogach configs oraz runs.

\section{Wpływ przepływności bitowej na jakość rekonstrukcji}
\label{sec:wplyw-przeplywnosci}

% TODO: Porównać warianty różniące się liczbą słowników, rozmiarem codebooków lub krokiem czasowym kodera.
% Należy zestawić nominalną przepływność z metrykami jakości oraz z obserwacjami odsłuchowymi.

\section{Wpływ parametrów modelu na skuteczność kompresji}
\label{sec:wplyw-parametrow-modelu}

% TODO: Przeanalizować znaczenie wybranych parametrów: wymiaru latentnego, rvq-code-dim, funkcji aktywacji,
% strat STFT/mel, stabilizacji RVQ, bloków dekodera i post-refinera.

\section{Ocena jakości odtworzonego dźwięku}
\label{sec:ocena-jakosci}

% TODO: Zaprezentować wyniki metryk SI-SDR, PESQ, STOI, ViSQOL, LSD, L1 i cos.
% W przypadku przykładów audio należy wskazać lokalizację plików w dodatkach lub repozytorium,
% zamiast wklejać obszerne dane binarne do pracy.

\section{Dyskusja wyników}
\label{sec:dyskusja-wynikow}

% TODO: Omówić kompromisy między przepływnością, zrozumiałością, naturalnością, stabilnością treningu
% i kosztem obliczeniowym. Należy wskazać, które wyniki potwierdzają realizację celu pracy,
% a które ujawniają ograniczenia metody.
